\hspace{3mm}

\centerline{\bf \Huge Процессы и операции}

\centerline{\bf \large Базовый уровень}

\begin{flushright}
        \scriptsize{}
\end{flushright}

\problem{0.1} 

\problem{0.2} 


\problem{1}
На столе лежат в ряд пять монет: средняя --- орлом вверх, а остальные --- решкой вверх. За одну операцию разрешается одновременно перевернуть ровно три монеты, лежащие рядом. Можно ли, выполнив такую операцию несколько раз, добиться того, чтобы все пять монет лежали орлом вверх?

\problem{2}
В некоторой стране из каждого города выходит нечетное число дорог. На центральной площади каждого города поднят черный или белый флаг. Каждое утро в одном из городов, у которого число соседей с флагами другого цвета больше половины, меняют цвет флага. Может ли этот процесс продолжаться бесконечно?

\problem{3}
Катя написала длинное число из нулей и единиц 10101010101010. Даша может дописать к числу Кати в любом месте 1010 или зачеркнуть кусочек 01. Сможет ли она такими операциями получить 01?

\problem{4}
На доске написаны числа 1, 2,..., 10. Каждую секунду выбираются какие-то два числа и оба заменяются на их полусумму. Может ли через некоторое время на доске оказаться 5 пятерок и 5 четверок?

\problem{5}
На доске написано число 1. С числом разрешается проводить две операции: умножать его на 2 и переставлять в нем цифры. Можно ли с помощью таких операций получить число 123456?



\newpage

\hspace{3mm}

\centerline{\bf \Huge Процессы и операции}

\centerline{\bf \large Продвинутый уровень}

\begin{flushright}
        \scriptsize{}
\end{flushright}

\problem{1} 
Каждый из 20 учеников ЭлМШ участвовавших на заключительном отборе в Сириус узнал свои баллы, известные только ему. За один телефонный разговор двое сообщают друг другу свои баллы. Докажите, что за 36 разговоров все они могут узнать баллы друг друга.


\problem{2} 
На доске написаны два числа. Каждый день Баир стирает с доски оба числа и пишет вместо них их среднее арифметическое и среднее гармоническое. Утром первого дня на доске были написаны числа 1 и 2. Найдите произведение чисел, записанных на доске вечером 2024-го дня.


\problem{3}
В клетки таблицы  $2024 \times 2024$ вписаны некоторые числа. Разрешается одновременно менять знак у всех чисел одного столбца или одной строки. Докажите, что несколькими такими операциями можно добиться того, чтобы суммы чисел, стоящих в любой строке и в любом столбце, были неотрицательны.


\problem{4}
На доске написаны числа 1, 2,..., 10. Каждую секунду выбираются какие-то два числа и оба заменяются на их полусумму. Может ли через некоторое время на доске оказаться 5 пятерок и 5 четверок?


\problem{5}
$\bigl[\textbf{Реальный ЕГЭ, 2023, последнее задание}\bigr]$
Есть числа A и B. Из них можно сделать числа $A + 2$ и $B - 1$ или $B + 2$ и $A - 1$, только если следующая пара этих чисел будет натуральной. Известно, что $A = 7, B = 11.$

а) Можно ли за 20 ходов создать пару, где одно из чисел равно 50?

б) За сколько ходов можно сделать пару, где сумма чисел будет равна 600?

в) Какое наибольшее число ходов можно сделать, чтобы оба числа не превышали 50?